\subsubsection{$M$ 理论预备}\label{prelims}

$\sf A_1$ 型的六维 $(0,2)$ 模型可以在一对平行的 M5 膜上实现。因此我们从坐标 $x^0,\dots,x^{10}$ 的 $\R^{11}$ 出发，考虑由 $x^0,\dots,x^5$ 参数化、位于 $x^6=\dots=x^{10}=0$ 的两张 M5 膜。该系统保持 16 个整体超对称。其生成元可理解为满足
\begin{equation}\label{zog}\Gamma_0\Gamma_1\cdots\Gamma_5\varepsilon=\Gamma_6\Gamma_7\cdots
 \Gamma_{10}\varepsilon=\varepsilon. \end{equation}
的十一维旋量 $\varepsilon$。这里 gamma 矩阵满足 $\{\Gamma_\mu,\Gamma_\nu\}=2g_{\mu\nu}$。

% [Tutor Note: 易错点] 这里 gamma 矩阵的约定是 $\{\Gamma_\mu,\Gamma_\nu\}=2g_{\mu\nu}$（含度规，而非 $2\delta_{\mu\nu}$）。因此 $\Gamma_\mu^2=g_{\mu\nu}$ 的符号依赖度规号差，$\Gamma_0\Gamma_1\cdots\Gamma_5$ 这类乘积的平方为 $\pm1$；后面 (\ref{omly}) 中 $\chi=\Gamma_0\Gamma_1$ 的平方为 $-1$（取决于号差），故 $\chi$ 的本征值为 $\pm i$ 还是 $\pm1$ 需与作者的号差约定一起核对。

为简化图像，我们引入对称破缺，把两张 M5 膜沿（比如）$x^6$ 方向分开。于是把一张置于 $x^6=0$，另一张置于 $x^6=L$，$L$ 为某个常数。在低能下，该系统由一对阿贝尔张量多重态描述，并与 BPS 弦耦合。这些弦来自张紧在两张 M5 膜之间的 M2 膜。弦的张力为 $T=T_{\mathrm{M2}}L$，其中 $T_{\mathrm{M2}}$ 是 M2 膜的张力。

我们将考虑一条世界体由 $x^0$ 与 $x^1$ 参数化、并位于指定 $x^2,\dots,x^5$ 处的弦。这条弦当然由一张从 $x^6=0$ 张紧到 $x^6=L$ 的 M2 膜表示，不过该方向将在下面的分析中直接分离出去。

上一段所描述的弦保持那些其生成元满足
\begin{equation}\label{zorf}\Gamma_0\Gamma_1\Gamma_6\varepsilon=
 \varepsilon\end{equation}
并且满足 (\ref{zog}) 的超对称。总共有八个未破缺超对称，对应于二维意义上的 $\N=4$ 超对称。有一个转动 $x^2,\dots, x^5$ 坐标的 $SO(4)$ 对称群，还有第二个 $SO(4)$ 对称，我们称之为 $SO(4)_R$，它转动 $x^7,\dots,x^{10}$。


然而，我们很快就会以某种方式修改构造，使一半超对称被破坏，并把 $SO(4)\times SO(4)_R$ 约化为一个极大环面。因此从一开始就着眼于相关对称性会有帮助。我们考虑由如下超对称组成的 $\N=2$ 子代数：它们（在已给条件之外）在（比如）$x^4-x^5$ 平面的转动与 $x^9-x^{10}$ 平面的 $R$ 对称转动的联合变换下不变：
\begin{equation}\label{dorz}
\left(\Gamma_4\Gamma_5+\Gamma_9\Gamma_{10}\right)\varepsilon=0.
\end{equation}
给定 (\ref{zorf})，这等价于
\begin{equation}\label{orz}\left(\Gamma_2\Gamma_3+\Gamma_7\Gamma_8\right)\varepsilon=0.
\end{equation}
我们注意到，若把 $W=x^2+ix^3$ 与 $Z=x^4+ix^5$ 互换，并相应地互换 $x^7,x^8$ 与 $x^9,x^{10}$，则方程 (\ref{dorz}) 与 (\ref{orz}) 互换。因此特别地，我们对 $\varepsilon$ 的条件在 $W$ 与 $Z$ 之间是对称的。

% [Tutor Note: 省略推导] 作者称“给定 (\ref{zorf})，(\ref{dorz}) 等价于 (\ref{orz})”，略去了代数步骤：先用 (\ref{zorf}) 把 $\Gamma_0\Gamma_1\Gamma_6$ 换成 1，再乘以 $\Gamma_2\Gamma_3\Gamma_4\Gamma_5$ 之类的组合，并利用 $\Gamma_4\Gamma_5$ 与 $\Gamma_9\Gamma_{10}$ 与其余因子的反对易性，即可把 $(\Gamma_4\Gamma_5+\Gamma_9\Gamma_{10})\varepsilon=0$ 化为 $(\Gamma_2\Gamma_3+\Gamma_7\Gamma_8)\varepsilon=0$。核对时请注意这一步用到了 $(\Gamma_4\Gamma_5)^2=-1$（在 $\{\Gamma_\mu,\Gamma_\nu\}=2g_{\mu\nu}$ 约定下）。

由上列条件挑出的 $\N=2$ 超对称代数具有一个 $U(1)\times U(1)$ 的 $R$ 对称群，其生成元为 $J=\Gamma_7\Gamma_8$ 与 $J'=\Gamma_9 \Gamma_{10}$。我们希望把 $J$ 与 $J'$ 和另外两个我们记作 $J_+$、$J_-$ 的 $R$ 对称生成元作比较，后者分别只作用于具有正、负二维手征性的超对称生成元，换言之，只作用于满足 $\chi\varepsilon=\pm \varepsilon$ 的生成元，其中 $\chi=\Gamma_0\Gamma_1$ 是二维手征性。上列条件可以合并给出
\begin{equation}\label{omly}\chi\varepsilon=JJ'\varepsilon\end{equation}
这意味着（对 $J_+$ 与 $J_-$ 作适当符号选择，并把它们归一化使得在相应手征性的超对称生成元上平方为 1）$J$ 与 $J'$ 可以表示为 $J=J_++J_-$，$J'=J_+-J_-$。

特别地，交换 $J\leftrightarrow J'$ 相当于 $J_\pm\to \pm J_\pm$，这一操作即所谓的 $\N=2$ 代数的镜像对称自同构。上述 $\N=2$ 代数中交换 $W$ 与 $Z$ 的自同构也交换 $J$ 与 $J'$，因此它是一个镜像对称。于是，若我们在弦的二维世界面理论中把 $W$ 视为手征超场，就必须按 \cite{GHR} 的意义把 $Z$ 视为扭手征超场。

\subsubsection{约化为规范理论}\label{gaugered}

到此为止我们只有半 BPS 弦，还没有它们的规范理论描述。

要得到规范理论描述，我们把一个方向——比如 $x^5$ 方向——紧致化在半径为 $R$ 的圆上。圆上的 $M$ 理论在长距离下约化为 IIA 型超弦理论。M5 膜变成 D4 膜，而 D4 膜上的理论在长距离下是 $U(2)$ 规范理论，并被 D4 膜之间的间隔破缺为 $U(1)\times U(1)$。与我们相关的是 $SU(2)$ 子群，它在低能下破缺为 $U(1)$。原先由张紧的 M2 膜所描述的弦，现在由张紧在两张 D4 膜之间的 D2 膜表示。

这一情形已在文献~\cite{diac} 中研究过。张紧在两张 D4 膜之间的 D2 膜带磁荷，对应于低能 $SU(2)$ 规范理论中的一个光滑 BPS 磁单极子。

在紧致化 $x^5$ 方向之前，沿弦的低能理论由手征超场 $W=x^2+ix^3$ 与扭手征超场 $Z=x^4+ix^5$ 描述。紧致化之后，我们可以把 $Z$ 换为单值场 $\exp(Z/R)$，它仍然是扭手征超场。然而，为了与 BPS 磁单极子的理论相匹配，另一个变量更为有用。对应于 D2 膜的 BPS 磁单极子的模是底层 M2 膜在 $x^2,x^3,$ 与 $x^4$ 中的位置，以及它在 $x^5$ 中位置的{\it 对偶}。这里所说的对偶是支配该弦的二维有效场论中的 $T$ 对偶。这个 $T$ 对偶是二维意义上的镜像对称。它把 $x^5/R$ 换成一个新的角度坐标 $\vartheta$。因此，$Z/R=x^4/R+ix^5/R$ 是扭手征超场，而 $Y=x^4/R+i\vartheta$ 是普通的手征超场，与 $W$ 一样。当然，单值的手征超场不是 $Y$ 而是 $e^Y$。

% [Tutor Note: 物理图像] $x^5/R$ 与 $\vartheta$ 的 $T$ 对偶把“沿 $x^5$ 的 Kaluza-Klein 动量”与“$U(1)$ 电荷/瞬子数”互换。这正是 M5 膜在圆上的理论与 D4 膜规范理论关系的一部分：转动 $x^5$ 的对称性在 $4+1$ 维规范理论中成为瞬子数，而转动 $\vartheta$ 的对称性成为电荷。也正因为 $\vartheta$ 是角度变量，$e^Y$（而非 $Y$）才是单值的，这与第~\ref{outline}~节 $\sY$ 的多值性讨论一致。

我们以两点评论作结：
\begin{itemize}
\item BPS 磁单极子的角度坐标 $\vartheta$ 与角位置 $x^5$ 互为 $T$ 对偶，这是“圆上的 M5 膜”与“D4 膜规范理论”之间关系的一部分。转动 $x^5$ 的对称性在 $4+1$ 维规范理论中成为瞬子数，而转动 $\vartheta$ 的对称性是电荷。
\item 单值场 $\Omega=\exp(Z/R)$ 可以理解为一个到 $\C^*$ 的映射；它既不能为零也不能为无穷大。在第~\ref{mod}~节我们将修改问题，使得 $\Omega=0$ 成为可能。
\end{itemize}

\subsubsection{在半空间上约化}\label{mod}
要得到非平凡的超势，我们必须破坏问题的一部分平移对称性。事实上，我们关心的是半空间上的规范理论，因此要把 $y=x^4$ 限制为非负。

本文所研究的规范理论问题，出现在 $x^4$ 与 $x^5$ 并不像我们迄今的表述中那样参数化 $\R\times S^1$，而是参数化一份带有雪茄型度规的 $\R^2$ 之时：
\begin{equation}\label{cigarl}\d s^2=\d y^2+f(y)\left(\d x^5\right)^2. \end{equation}
其中 $f(y)\sim y^2/R^2$（$y\to 0$），且 $f(y)\to 1$（$y\to\infty$）。事实上，这正是文献~\cite{fiveknots} 中推导的出发点。

作为复流形，$\R^2$ 与 $\C$ 相同，因此现在可以用一个取值于 $\C$ 的手征超场 $\Omega$ 来参数化 $x^4$ 与 $x^5$ 方向；在大 $y$ 处（且仅在那里）渐近地有
\begin{equation}\label{exo}\Omega=\exp(Z/R), ~~Z=x^4+ix^5.\end{equation}
与第~\ref{gaugered}~节最后的评论不同，$\Omega$ 现在取值于 $\C$ 而非 $\C^*$，特别地，取 $\Omega=0$ 不再有任何问题。

我们很快就会在生成超势时用到这一点。

\subsubsection{瞬子}\label{thinst}

现在我们要描述一个将生成我们所寻找的超势的 M2 膜瞬子。

该瞬子应当修正一条由世界面坐标 $X=x^0+ix^1$ 参数化的弦的物理。这条弦位于 $W$ 与 $\Omega$ 的确定值处，比如 $W=W_0$、$\Omega=\Omega_0$。

我们可以定性地理解何种瞬子能够生成超势。考虑一个具有 $\N=2$ 超对称的二维模型，其手征环在经典层面由手征超场 $W$ 生成，其扭手征环由扭手征超场 $\Omega$ 生成。手征环是扭 $B$ 模型的可观测量构成的环，而我们要生成的超势将给出该手征环的一个形变。一个假想的超势必须由保持 $B$ 模型超对称的位形生成。这类位形中，手征超场 $W$ 为常数，而扭手征超场 $\Omega$ 作为世界面坐标的函数是全纯的。（众所周知，$B$ 模型超对称要求手征超场为常数。$B$ 模型超对称允许扭手征超场全纯，这一点与或许更为人熟知的“$A$ 模型超对称允许手征超场全纯”互为镜像。）

因此，生成超势的瞬子由 $W=W_0$ 给出，而 $\Omega$ 是 $X$ 的一个非平凡但简单的全纯函数：
\begin{equation}\label{yelf}\frac{\Omega}{\Omega_0}=X-X_0,\end{equation}
（注意，这个公式之所以有意义，正是因为允许 $\Omega$ 取零。）这里 $X_0$ 是一个常数，我们把它解释为瞬子的位置；正如瞬子物理中一贯的那样，为计算物理振幅，必须对瞬子模积分，这里即指对 $X_0$ 以及某些与瞬子解在其下不变的那些超荷相联系的费米模积分。

M2 膜瞬子的世界体应当是一个三维流形。我们想要的三维流形正是 (\ref{yelf}) 所定义的二维流形 $S$ 与一维流形 $0\leq x^6\leq L$ 的乘积（全部位于 $x^7=\dots=x^{10}=0$）。

由于它保持 $B$ 模型超对称，并且除了平移不变性与超对称所隐含的模之外没有其他模（参数 $\Omega_0$ 大致对应于没有瞬子时 $\Omega$ 会取的常数值），这类瞬子将生成一个超势。为理解究竟会生成怎样的超势，我们利用渐近公式 (\ref{exo})，考察瞬子在远距离处对弦产生的扰动。对大的 $y=x^4$，可以写出
\begin{equation}\label{omo}\exp((y+ix^5)/R)=\Omega_0(X-X_0).\end{equation}
我们看到，当 $X$ 沿顺时针方向绕 $X_0$ 一圈时（在 $|X-X_0|$ 很大、从而使公式 (\ref{omo}) 成立的区域），$x^5$ 增加 $2\pi R$。为产生这一效应，插入在 $X=X_0$ 处的算符必须是一个 twist field。由于 $x^5$ 的 $T$ 对偶是角度变量 $\vartheta$，一个 twist field 就是 $\exp(-i\vartheta)$，而这必定是刻画瞬子效应的超势对 $\vartheta$ 的依赖。因此全纯的表达式必须是 $\W=\Lambda \exp(-Y)$，其中 $Y=y+i\vartheta$，$\Lambda$ 为某个常数。这正是 (\ref{forg}) 所断言的结果。

% [Tutor Note: 物理图像] (\ref{omo}) 的要点是：把 $X\mapsto X-X_0$ 的对数单值性翻译成 $x^5\to x^5+2\pi R$，于是插入算符必须是“绕数”算符 $e^{-i\vartheta}$。$-i$（而非 $+i$）的符号由 (\ref{omo}) 中 $X$ 的顺时针绕行方向决定；若绕行方向或 $X=x^0+ix^1$ 的复结构约定相反，$\W=\Lambda\exp(-Y)$ 中会出现 $\exp(+Y)$，从而导致超势随 $y$ 指数增长而非衰减。

为进一步确认，并检验 $\W$ 指数上的符号，我们来考察场 $y$ 在远离 $X=X_0$ 的远距离处的行为。在长距离下，$y$ 的涨落由自由场路径积分描述
\begin{equation}\label{donel}\int Dy\,\exp\left(-\frac{1}{4\pi R^2}\int \d x^0\,\d x^1\,|\nabla y|^2\right).
\end{equation}
当算符 $\exp(-y/R)$ 在 $X=X_0$ 处插入这样的路径积分时，其结果是：在远距离处 $y/R$ 按 $|\log(X-X_0)|$ 增长。而这正是我们在 (\ref{omo}) 中看到的。

最后一点评论：如果弦的世界体维数大于 2，我们就会把瞬子看作围绕一个真空的涨落，该真空由 $y$（以及所有其他世界体场）在 $X\to\infty$ 时的极限值定义，并要求瞬子在无穷远处趋于该极限值。在二维情形，由于通常的红外发散——例如上一段提到的对数增长——这样的表述并不成立。




\section{oper 与膜}\label{opbranes}

\begin{keytakeaways}
\begin{itemize}\setlength{\itemsep}{2pt}
  \item \textbf{物理图像}：把四维 $\N=4$ 超 Yang-Mills 理论放在 $M_4=\R\times C\times\R_+$ 上并沿 $C$ 紧致化，低能得到靶空间为 Hitchin 模空间 $\M_H$ 的二维 sigma 模型；$\R_+$ 有限端的 Nahm 极点边界条件在低能下应约化为 $\M_H$ 中的一个膜，且因原边界条件为半 BPS 而该膜也必须是半 BPS 的。
  \item \textbf{数学工具}：Hitchin 纤维化（$(E,\varphi)\mapsto\Tr\varphi^2$）、$\M_H$ 上的三个复结构 $I$（Higgs 丛）、$J$（平坦丛 $\A=A+i\phi$）、$K=IJ$、复辛形式 $\Omega_I$、oper（满足非分裂正合列 $0\to\L\to E\to\L^{-1}\to0$ 的平坦丛）。
  \item \textbf{主要结论}：$\t=1$ 时 Nahm 极点对应的膜支撑在 Hitchin 纤维化的全纯截面 $\T$ 上，是 $(B,A,A)$ 型膜（对 $I$ 为 $B$ 膜、对 $J,K$ 为 $A$ 膜），$\T$ 对 $I$ 是复 Lagrangian 的；一般 $\t$ 时支撑在 opers 构成的 $\V_\zeta$ 上，$\V_\zeta$ 对旋转复结构 $I_\zeta$ 是复 Lagrangian 的，故为 $(B,A,A)_\zeta$ 型膜。
  \item \textbf{与本文的联系}：$\zeta\to0$ 时 $\V_\zeta$ 退化为 $\T$；$\zeta=i$ 时给出关于 $J$ 的通常 opers 簇，这正是 $S$ 对偶（第~\ref{sduality}~节）中 oper 膜 $\Boper$ 的几何基础。
\end{itemize}
\end{keytakeaways}

本节的目的在于把本文中出现的一些要素置于更宽广的背景之中。

我们继续研究四维流形 $M_4=\R\times C\times \R_+$ 上带保持一半超对称的 twist 的 $\N=4$ 超 Yang-Mills 理论，并在 $\R_+$ 的有限端施加通常的 Nahm 极点边界条件。与迄今所述相比，新颖之处在于我们将从沿 $C$ 作四维到二维紧致化的视角来看待该问题。一般而言~\cite{BVJ,HM}，为简单起见假定 $C$ 的亏格至少为 2（我们在第~\ref{mondr}~节放宽这一条件），沿 $C$ 的紧致化在低能下给出一个二维 sigma 模型，其靶空间是 $\M_H$，即 Hitchin 方程解构成的模空间~\cite{Hitchin}。Nahm 极点边界条件在低能下必须约化为该 sigma 模型中的一个膜，而且这个膜必须是半 BPS 的，因为 Nahm 极点边界条件在四维中是半 BPS 的。



\subsection{回到 $\t=1$}\label{tone}

我们首先分析 $\t=1$ 的情形。为简单起见，取 $G$ 为 $SU(2)$ 或 $SO(3)$。在第~\ref{bc}~节我们看到，在 $\t=1$ 时，Nahm 极点边界条件（在没有奇异磁单极子时）描述一个希格斯丛 $(E,\varphi)\to C$，它带有一个处处不被 $\varphi$ 保持的线子丛 $\L\subset E$。把 $E$ 看作行列式平凡的二阶复丛，则包含映射 $\L\subset E$ 是一个正合列的一部分：
\begin{equation}\label{exacto} 0 \to \L\to E\to \L^{-1}\to 0.\end{equation}
这里我们用到了如下事实：由于 $E$ 的行列式平凡，商丛 $E/\L$ 必然同构于 $\L^{-1}$。

我们把 $\varphi$ 视为全纯映射 $E\to E\otimes K$，其中 $K$ 是 $C$ 的典范丛。可以把 $\varphi$ 限制到 $\L$ 上，得到全纯映射 $\L\to E\otimes K$，再利用投影 $E\to \L^{-1}$，得到映射 $\varphi:\L\to \L^{-1}\otimes K$。$\L$ 处处不被 $\varphi$ 保持这一条件恰恰意味着映射 $\varphi:\L\to \L^{-1}\otimes K$ 处处非零。换言之，这个映射是同构。

与 $\L$ 作张量积，我们得知 $\L^2$ 同构于 $K$，因此 $\L$ 是 $K$ 的一个平方根 $K^{1/2}$。若 $G=SU(2)$，Nahm 极点边界条件的一个解涉及对 $K^{1/2}$ 的选取；而若 $G=SO(3)$，由于我们真正应当处理的是伴随丛 $\mathrm{ad}(E)$ 而非 $E$，$K^{1/2}$ 的选取无关紧要。以下我们假定或者 $G=SO(3)$，或者已经选定了 $K$ 的一个特定平方根。

可以构造非平凡扩张 $0\to K^{1/2}\to E\to K^{-1/2}\to 0$，我们将在第~\ref{general}~节利用这一事实。然而，对于希格斯丛，我们希望 $E$ 是直和 $K^{1/2}\oplus K^{-1/2}$，因为若扩张非平凡，我们即将写出的希格斯场将不存在。若把 $E$ 写成列向量形式
\begin{equation}\label{equad} E=\begin{pmatrix}K^{-1/2}\cr K^{1/2}\end{pmatrix}.\end{equation}
则在相差 $E$ 的自同构的意义下，一个可能的希格斯场 $\varphi$ 形如
\begin{equation}\label{varph}\varphi=\begin{pmatrix}0 & 1\cr q & 0 \end{pmatrix},\end{equation}
其中 $q$ 是一个二次微分。更准确地说，我们假定 $\varphi$ 的右上矩阵元非零，否则 $\L$ 就会被 $\varphi$ 保持（而正如文献~\cite{Hitchin} 所解释的，希格斯丛 $(E,\varphi)$ 将是不稳定的）。在此前提下，通过丛自同构 $\mathrm{diag}(\lambda,\lambda^{-1})$ 可把右上矩阵元取为 1，再通过下三角的丛自同构可令 $\varphi$ 的对角矩阵元为零。最后，对如 (\ref{equad}) 中的 $E$，$\varphi$ 的左下矩阵元是一个二次微分（即 $H^0(C,K^2)$ 中的元素），我们记之为 $q$。

% [Tutor Note: 省略证明] 作者称“通过下三角丛自同构可令对角矩阵元为零”，这一步需注意：下三角自同构 $\begin{pmatrix}1&0\\ \epsilon&1\end{pmatrix}$ 作用在 $\varphi$ 上会把 $q$ 换成 $q+c$ 而保持右上元为 1，并同时改变对角元；只有当对角元之差可由 $\epsilon$ 消去时才可行，这依赖于右上元非零（从而 $E$ 分裂为 $\L\oplus\L^{-1}$ 的直和）。若扩张非平凡，则 $H^0$ 中不存在这样的整体 $\epsilon$，这正是作者强调“对希格斯丛要求直和”的原因。

设 $\T\subset \M_H$ 是参数化上一段所述希格斯丛 $(E,\varphi)$ 的子流形。在 $\t=1$ 时，由 Nahm 极点定义的 $\M_H$ 中的膜支撑在 $\T$ 上。$\T$ 是怎样的子簇？如文献~\cite{Hitchin} 那样，设 $I$ 是 $\M_H$ 上使其参数化希格斯丛的复结构，$J$ 是使 $\M_H$ 参数化带联络 $\A=A+i\phi$ 的平坦丛的复结构，$K=IJ$。Hitchin 纤维化是把 $\M_H$ 映到二次微分空间的映射，它把 $(E,\varphi)$ 映为 $\mathrm{Tr}\,\varphi^2$。该映射在复结构 $I$ 下是全纯的。对 (\ref{equad}) 中的希格斯场，有 $\Tr\,\varphi^2=2q$，因此对每个给定的 $\Tr\,\varphi^2$ 值都存在唯一的这种 $\varphi$。相应地，$\T$ 是 Hitchin 纤维化的一个全纯截面；事实上它就是文献~\cite{Hitchin} 中所构造的全纯截面。实际上，从复结构 $I$ 的角度看，$\T$ 是复 Lagrangian 的。该断言意味着复辛形式
\begin{equation}\label{Om}\Omega_I=\frac{1}{4\pi}\int_C\d\bar z\,\d z\,\Tr\,\delta A_{\bar z}\delta\phi_z
\end{equation}
限制到 $\T$ 上时为零。事实确实如此，因为对代表 $\T$ 中点的所有希格斯丛，$E$ 的全纯类型都是固定的，从而限制到 $\T$ 上时 $\delta A_{\bar z}$ 为零（相差一个规范变换）。

由于 $\T$ 在复结构 $I$ 下是复 Lagrangian 的，我们可以如下识别 Nahm 极点所产生的半 BPS 膜的超对称类型。这是 $(B,A,A)$ 型膜，即它在复结构 $I$ 下是 $B$ 膜，而从 $J$ 或 $K$ 的角度看是 $A$ 膜。

\subsection{一般的 $\t$}\label{general}

在一般的 $\t$ 下，我们处理的是平坦丛而非希格斯丛。Nahm 极点仍然给出线子丛 $L\subset E$，因此我们仍有正合列
\begin{equation}\label{texaco} 0\to \\ E\to \L^{-1}\to 0,\end{equation}
与 (\ref{exacto}) 一样。协变导数 $\D/\D z$ 现在给出全纯映射 $E\to E\otimes K$。我们仍可把这个映射限制到 $\L$ 上，并把像投影到 $\L^{-1}\otimes K$，得到线性映射 $\D/\D z:\L\to \L^{-1}\otimes K$。$\L$ 处处不被 $\D/\D z$ 保持这一条件，正如第~\ref{tone}~节那样，意味着该映射是从 $\L$ 到 $\L^{-1}\otimes K$ 的同构，于是我们再次得出结论：$\L=K^{1/2}$。

% [Tutor Note: 原文笔误] (\ref{texaco}) 中写作 $0\to \\ E$，其中出现了双反斜杠（原文如此），LaTeX 会把它解释为换行符而非丛记号 $\L$。此处公式按原文逐字保留，未作修改；编译时若报错，应核对是否应为 $0\to \L\to E$。此外本段第一句“Nahm 极点仍给出线子丛 $L\subset E$”用了 $L$ 而非 $\L$，与全文记号不一致，亦按原文保留。

与第~\ref{tone}~节的差别在于，现在丛 $E$ 不是直和 $K^{1/2}\oplus K^{-1/2}$，而是非平凡扩张。确实，由于我们假定 $C$ 的亏格大于 1，一个在全纯意义上为直和 $K^{1/2}\oplus K^{-1/2}$ 的丛不可能容许平坦联络。

$K^{1/2}$ 被 $K^{-1/2}$ 的非平凡扩张彼此都同构；这是因为这样的扩张由 $H^1(C,K)\cong \C$ 中的一个元素确定，而在相差丛自同构的意义下，非零元素的选取无关紧要。


从全纯观点看是上一段所述扩张的平坦丛，其最简单的例子可以通过在 $C$ 上放置标量曲率 $R=-1$ 的 Kähler 度规来得到。设 $\omega$ 为该度规的自旋联络，$e$ 为四元标架（vierbein）。平坦联络为
\begin{equation}\label{omurf}\A=\omega\ct_3+e_{\bar z}\ct_- +e_z\ct_+.\end{equation}
在微分几何中，由于 $\omega$ 是自旋联络，平坦丛 $E$ 就是 $C$ 的自旋丛，更准确地说，是相反手征性的两个自旋丛的直和 $K^{1/2} \oplus K^{-1/2}$。但在这个基下，$E$ 的复结构由 $\A$ 的 $(0,1)$ 部分定义，即 $\A_{\bar z}=\omega_{\bar z}\ct_3 +e_{\bar z}\ct_-$；它是下三角的但不是对角的，因此 $E$ 是扩张而非直和。

在丛 $E$ 上找到一个平坦联络 $\A$ 之后，找出全部平坦联络是容易的。我们不想改变 $\A_{\bar z}$（因为 $E$ 的全纯结构应保持不变），但可以把 $\A_z$ 改为 $\A_z\to \A_z+\lambda_z$，其中（为保持平坦性）$\lambda_z$ 被 $\D_{\bar z}$ 湮灭。对上一段所述的 $E$，相关的选择是 $\lambda_z=q\ct_-$，其中 $q$ 是一个二次微分。

在数学上，从全纯观点看适合非分裂正合列 (\ref{texaco}) 的平坦丛 $E\to C$ 称为 oper；详见文献~\cite{Frenkel}。我们已经知道，在一般的 $\t$ 下，由 Nahm 极点边界条件定义的膜支撑在 opers 构成的簇上。实际上我们应当更精确些，因为 Nahm 极点边界条件一般依赖于第~\ref{motivation}~节引入的参数 $\zeta$，而满足 oper 条件的复联络 $\A$ 一般不是 $A+i\phi$，而是式 (\ref{generic}) 中所定义的更一般的联络 $\A^\zeta$。我们把 $\M_H$ 中要求 $\A^\zeta$ 满足 oper 条件的子簇记作 $\V_\zeta$。



限制到 $\V_\zeta$ 上时，$\A^\zeta_{\bar z}$ 在相差规范变换的意义下是固定的，因此复辛形式
\begin{equation}\label{tufo}\Omega_{I_\zeta}=\frac{1}{4\pi}\int_C\d\bar z\,\d z\,\Tr\,\delta\A^\zeta_{\bar z}\delta\A_z
^\zeta
\end{equation}
为零。相应地，膜 $\V_\zeta$ 是复 Lagrangian 膜，正如第~\ref{tone}~节那样，只不过现在是在旋转后的复结构中。就本文而言，旋转后的复结构是第~\ref{motivation}~节定义的 $I_\zeta$。$\V_\zeta$ 也许可以称为 $(B,A,A)_\zeta$ 型膜：它与复结构 $I_\zeta$ 的关系，正如 $(B,A,A)$ 膜与复结构 $I$ 的关系。

只要 $\zeta\not=0,\infty$，各复结构 $I_\zeta$ 都是等价的。若简单地取 $\zeta=i$，就得到关于复结构 $J$ 的通常的 opers 簇；另一方面，在 $\zeta\to 0$ 的极限下，$\V_\zeta$ 退化为第~\ref{tone}~节所述的 Hitchin 纤维化的全纯截面 $\V$。

% [Tutor Note: 约定差异] 复辛形式 (\ref{Om}) 与 (\ref{tufo}) 的归一化因子 $1/4\pi$ 与符号约定在各文献中并不统一；本文的 $\Omega_I$ 取 $\delta A_{\bar z}\,\delta\phi_z$ 的缩并，而 $\Omega_{I_\zeta}$ 取 $\delta\A^\zeta_{\bar z}\delta\A^\zeta_z$。判断“复 Lagrangian”时只需 $\Omega$ 限制为零，与该归一化无关，但若要比较不同文献给出的 $\Omega$ 的整数系数（例如是否多出 $1/2$），必须回到这两个定义。
